\startslide
\commonframe
\mbox{}

\begin{tikzpicture}[remember picture,overlay]
  \slidetitle{Модель данных}

  \node[anchor=north west, align=left, text width=6.0cm] at ([xshift=1.00cm,yshift=-1.95cm]current page.north west) {%
    {\fontsize{9.9}{11.5}\selectfont
    Центральная сущность исследования --- \textbf{CRDT-операция}. Во всех трёх
    вариантах состав данных одинаковый, различается только физическая
    организация хранения.}%
  };

  \node[
    anchor=north west,
    fill=softgray,
    rounded corners=10pt,
    text width=6.0cm,
    inner xsep=0.35cm,
    inner ysep=0.28cm,
    align=left
  ] at ([xshift=1.00cm,yshift=-4.68cm]current page.north west) {%
    {\fontsize{10.0}{11.6}\selectfont
    \texttt{opId}\\
    \texttt{txId}\\
    \texttt{objectId}\\
    \texttt{replicaId}\\
    \texttt{clock}\\
    \texttt{action}}%
  };

  \node[
    anchor=north west,
    fill=softcyan,
    rounded corners=10pt,
    text width=6.95cm,
    inner xsep=0.38cm,
    inner ysep=0.24cm,
    align=left
  ] at ([xshift=7.65cm,yshift=-1.95cm]current page.north west) {%
    {\fontsize{9.1}{10.4}\selectfont
    \textbf{Смысл полей}\\[0.10cm]
    \texttt{clock} хранит причинно-следственную информацию,\\
    \texttt{action} хранит полезную нагрузку операции.\\[0.14cm]
    Поэтому работа сравнивает не разные данные, а разные
    \textbf{модели представления одного operation log}.}%
  };

  \node[anchor=north west, text=textblue] at ([xshift=7.65cm,yshift=-5.58cm]current page.north west) {%
    {\fontsize{10.2}{11.4}\selectfont\bfseries Единицы хранения}%
  };

  \node[anchor=north west, align=left, text width=6.95cm] at ([xshift=7.70cm,yshift=-5.96cm]current page.north west) {%
    {\fontsize{9.0}{10.0}\selectfont
    \textbf{PostgreSQL:} отдельная операция как строка таблицы.\\
    \textbf{MongoDB:} объект как документ со встроенной историей.\\
    \textbf{YTsaurus:} операция как строка distributed dynamic table.}%
  };
\end{tikzpicture}

\footerband{3}
